\section{طراحی مدار محاسباتی}

\newimage[\textwidth]{images/comb-circuit/comb-circuit.png}{مدار محاسباتی}{fig:comb-circuit}

ساختار فوق را با استفاده از کتابخانه‌ی \texttt{basic\_gates.inc} که از پروژه اول تا الان بروزرسانی شده است و عملکرد درستی دارند پیاده‌سازی می‌کنیم.
گیت‌ها به گونه‌ای پیاده سازی شده‌اند که با پارامتر scale می‌توان بدون درگیر شدن در محاسبات Sizing ابعاد آن‌ها را تغییر بدهیم.

\newcode[1]{SPICE}{codes/params.sp}{پارامتر‌های تعریف شده برای شبیه‌سازی}{code:param}
\newcode[1]{SPICE}{codes/samp-gate.sp}{نمونه گیت \lr{NOR3} از کتابخانه گیت‌های پایه}{code:sample-gate}

\begin{multicols}{2}
    \newimage[0.45\textwidth]{./images/gate/NAND2.png}{شماتیک گیت \lr{NAND2}}{fig:NAND2}
    \newimage[0.45\textwidth]{./images/gate/NOR2.png}{شماتیک گیت \lr{NOR2}}{fig:NOR2}
\end{multicols}

\begin{multicols}{2}
    \newimage[0.53\textwidth]{./images/gate/NAND3.png}{شماتیک گیت \lr{NAND3}}{fig:NAND3}
    \newimage[0.45\textwidth]{./images/gate/NOR3.png}{شماتیک گیت \lr{NOR3}}{fig:NOR3}
\end{multicols}

\begin{multicols}{2}
    \newimage[0.45\textwidth]{./images/gate/XOR2.png}{شماتیک گیت \lr{XOR2}}{fig:XOR2}
    \newimage[0.45\textwidth]{./images/gate/XNOR2.png}{شماتیک گیت \lr{XNOR2}}{fig:XNOR2}
\end{multicols}


\subsection{شبیه‌سازی مدار}

\newcode[1]{SPICE}{codes/comb-circuit.sp}{پیاده‌سازی \lr{HSPICE} مدار محاسباتی}{code:comb-circuit}

برای تولید \lr{transition‌}های شبیه‌سازی از اسکریپت پایتون \texttt{delay-sim.py} استفاده شده.
از اسکریپت‌ پایتون \texttt{modify-mtx.py} برای استخراج دقیق تاخیرها از فایل‌های خروجی شبیه‌سازی استفاده می‌کنیم. این اسکریپت‌ها به صورت خودکار \lr{Transition}های داده‌های ورودی را ساخته و داده‌های مربوط به تأخیر انتشار  را از فایل‌های \texttt{.mtx} استخراج کرده و در قالب فایل \texttt{.json} ذخیره می‌کنند.
\pfnote{این اسکریپت‌ها توسط \lr{LLM-Tools} تولید شده اند ولی هزار بار چک شده اند :D}
با این روش در هر idex از شبیه‌سازی ما یکی از حالات transition را بررسی می‌کنیم و مقدار تاخیر را اندازه‌گیری می‌کنیم.
در تمام شبیه‌سازی های مدارات محاسباتی در تمام بخش‌ها \lr{$t_{rise} = t_{fall} = 50ps$}  و برای مدارات ترتیبی \lr{20ps} برای سیگنال‌های ورودی اعمال شده است.
\newcodeFont[\tiny]{spice}{codes/delay-sim.sp}{بخشی‌ از .DATA ساخته شده توسط اسکریپت}{code:delay-sim}

\newimage[\textwidth]{images/comb-circuit/delay-sim.png}{شکل موج خروجی شبیه‌سازی‌ها به ازای تمام \lr{Transition}ها}{fig:delay-sim}

خروجی شبیه‌سازی مطابق انتظار است چون خروجی همیشه باید صفر باشد. مگر اینکه با یک glich حاصل از دو مسیر اول و دوم برای مدتی یک شود.
به این ترتیب تاخیر انتشار این مدار برابر با زمانی است که از تغییر ورودی تا stable شدن خروجی می‌گذرد.

\newimage[\textwidth]{images/comb-circuit/comb-circuit-2.png}{مدار محاسباتی}{fig:comb-circuit-2}

برای یک بودن خروجی می‌بایست تمام ورودی‌های آخرین گیت (که یک گیت \lr{NOR3} هست) صفر باشند.
ولی اگر خروجی \lr{XOR2} صفر شود، خروجی گیت \lr{NAND2}ای که drive می‌کند یک می‌شود!
اما در این میان برای مدتی به اندازه تاخیر انتشار گیت \lr{NAND2} هرسه ورودی می‌توانند صفر باشند.
ازونجایی که \lr{$V_{dd}/2$} روی ورودی‌ها در \lr{100ps} اتفاق می‌افتد می‌توان تأخیرها را اینگونه فرض کرد: \lr{$t_{cd}=250ps$} و \lr{$t_{pd}=400ps$}.

\newpage
مدار محاسباتی را تغییر می‌دهیم و به‌جای \lr{NAND2} با \lr{scale=2.4} یک \lr{AND2} با همین scale قرار می‌دهیم:

\newimage[0.78\textwidth]{images/comb-circuit/comb-circuit-3.png}{مدار محاسباتی}{fig:comb-circuit-3}

\newimage[0.85\textwidth]{images/comb-circuit/delay-sim2.png}{شکل موج خروجی شبیه‌سازی‌ها به ازای تمام \lr{Transition}ها}{fig:delay-sim2}

در نهایت این مقادیر تاخیر برای مدار محاسباتی measure شده است:

\lr{
    $$
        \begin{cases}
            t_{pd} = 541 \text{ ps}   & (5.41 \times 10^{-10} \text{ s})  \\
            t_{cd} = 424.2 \text{ ps} & (4.242 \times 10^{-10} \text{ s})
        \end{cases}
    $$
}

اسکریپت‌های استخراج داده به گونه‌ای طراحی شده‌اند که \lr{transition}ای که با آن این تاخیر‌ها بدست آمده‌اند را نیز ذخیره می‌کنند.

\newcodeFont[\scriptsize]{spice}{codes/comb-delay.sp}{\lr{Transition}هایی که تاخیر‌های اکسترمم را بدست‌ آورده‌اند}{code:comb-delay}
